% ---- pg-064 （书页 52） ----

\begin{mvtable}
\begin{tabular}{lccc}
\toprule
$x_1$ & $-1$ & $0$ & $1$ \\
\midrule
$p_1(x_1)$ & $0.3$ & $0.3$ & $0.4$ \\
\bottomrule
\end{tabular}
\end{mvtable}

则
\[
E(X_1)=\sum_{\text{一切 }x_1}x_1p_1(x_1)=(-1)\times0.3+0\times0.3+1\times0.4=0.1
\]

类似地，设离散型随机变量 $X_2$ 有概率函数

\begin{mvtable}
\begin{tabular}{lcc}
\toprule
$x_2$ & $0$ & $1$ \\
\midrule
$p_2(x_2)$ & $0.8$ & $0.2$ \\
\bottomrule
\end{tabular}
\end{mvtable}

则
\[
E(X_2)=\sum_{\text{一切 }x_2}x_2p_2(x_2)=0\times0.8+1\times0.2=0.2
\]

因此
\[
E(\bm{X})=
\begin{bmatrix}
E(X_1)\\
E(X_2)
\end{bmatrix}
=
\begin{bmatrix}
0.1\\
0.2
\end{bmatrix}
\]

直接由随机矩阵期望值的定义以及期望的性质 $E(X_i+Y_i)=E(X_i)+E(Y_i)$ 和 $E(cX_i)=cE(X_i)$，可得到关于矩阵的和的期望与乘积的期望的两个结论.设 $\bm{X}$ 和 $\bm{Y}$ 是同维的随机矩阵，并令 $\bm{A}$ 和 $\bm{B}$ 是常数的可相乘矩阵，则(参见练习 2.40)

\begin{mvframe}
\phantomsection\label{eq:2-24}
\begin{equation}
\begin{aligned}
E(\bm{X}+\bm{Y})&=E(\bm{X})+E(\bm{Y})\\
E(\bm{A}\bm{X}\bm{B})&=\bm{A}E(\bm{X})\bm{B}
\end{aligned}
\tag{2-24}
\end{equation}
\end{mvframe}

\section{均值向量和协方差矩阵}  % 原书 2.6
\phantomsection\label{sec:2.6}

设 $\bm{X}=[X_1,X_2,\cdots,X_p]$ 是 $p\times1$ 随机向量，则 $\bm{X}$ 的每个元素都是一个有着自己的边缘概率函数的随机变量(参考例 2.12).边缘均值 $\mu_i$ 和方差 $\sigma_i^2$ 分别定义为 $\mu_i=E(X_i)$ 和 $\sigma_i^2=E(X_i-\mu_i)^2,i=1,2,\cdots,p$.具体地说，

\[
\mu_i=
\begin{cases}
\displaystyle\int_{-\infty}^{\infty}x_if_i(x_i)\,\mathrm{d}x_i & \text{若 } X_i \text{ 是概率密度函数为 } f_i(x_i) \text{ 的连续型随机变量}\\[2ex]
\displaystyle\sum_{\text{一切 }x_i}x_ip_i(x_i) & \text{若 } X_i \text{ 是概率函数为 } p_i(x_i) \text{ 的离散型随机变量}
\end{cases}
\]

\[
\sigma_i^2=
\begin{cases}
\displaystyle\int_{-\infty}^{\infty}(x_i-\mu_i)^2f_i(x_i)\,\mathrm{d}x_i & \text{若 } X_i \text{ 是概率密度函数为 } f_i(x_i) \text{ 的连续型随机变量}\\[2ex]
\displaystyle\sum_{\text{一切 }x_i}(x_i-\mu_i)^2p_i(x_i) & \text{若 } X_i \text{ 是概率函数为 } p_i(x_i) \text{ 的离散型随机变量}
\end{cases}
\]

\phantomsection\label{eq:2-25}
\begin{equation}
\tag{2-25}
\end{equation}

在以后的各节中，采用记号 $\sigma_{ii}$ 而不是更传统的 $\sigma_i^2$ 来表示边缘方差，较为方便.因此，我们将使用这一记号.

任何一对随机变量 $X_i$ 和 $X_k$ 的特性，可用它们的联合概率函数来描述.它们之间的线性联系的程度可用协方差来表示.

\[
\sigma_{ik}=E(X_i-\mu_i)(X_k-\mu_k)
\]
